\section{Техническое задание}

\subsection{Основание для разработки}

Данная разработка ведётся на основании задания на выпускную квалификационную работу бакалавра по теме «Разработка модуля взаимодействия с пользователем».

\subsection{Цель и назначение разработки}

Целью разработки является создание серверного модуля взаимодействия с пользователем для онлайн площадки по аренде вещей. Модуль обеспечивает коммуникацию между участниками сделки, своевременное информирование пользователей о событиях платформы и автоматизированное формирование юридически значимой документации.

Задачами данной разработки являются:
\begin{itemize}
    \item проектирование и реализация подсистемы чатов с поддержкой передачи изображений и доставки сообщений в режиме реального времени через WebSocket;
    \item проектирование и реализация подсистемы уведомлений, включающей всплывающие оповещения на сайте через WebSocket и отправку email-сообщений при ключевых событиях платформы;
    \item проектирование и реализация подсистемы генерации договоров аренды в формате PDF с поддержкой демонстрационной электронной подписи обеих сторон сделки;
    \item разработка справочного раздела часто задаваемых вопросов (FAQ) по основным действиям в системе.
\end{itemize}

\subsection{Требования к модулю взаимодействия с пользователем}

\subsubsection{Функциональные требования}

На основании анализа требований к платформе сформирован перечень функциональных требований к разрабатываемому модулю, представленный в таблице~\ref{tz:table}.

\begin{xltabular}{\textwidth}{|c|X|X|}
 \caption{Функциональные требования к модулю взаимодействия с пользователем \label{tz:table}}\\ \hline
  № & \centrow Функциональное требование & \centrow Описание \\ \hline
 \endfirsthead
 \continuecaption{Продолжение таблицы \ref{tz:table}}
 ~ № & \centrow Функциональное требование & \centrow Описание \\ \hline 
 \finishhead
 \leftrow 1 & Создание и получение чатов & \leftrow Создание диалога между двумя пользователями, привязка чата к объявлению; получение списка чатов текущего пользователя, упорядоченных по дате последнего сообщения. \\ \hline
 \leftrow 2 & Отправка сообщений & \leftrow Отправка текстовых сообщений и изображений через REST API; немедленная доставка сообщения всем участникам чата через WebSocket без перезагрузки страницы. \\ \hline
 \leftrow 3 & Отметка сообщений как прочитанных & \leftrow Отметка непрочитанных сообщений чата как прочитанных при открытии диалога; рассылка события об изменении статуса всем участникам чата через WebSocket. \\ \hline
 \leftrow 4 & Счётчик непрочитанных сообщений & \leftrow Подсчёт и передача клиенту количества непрочитанных сообщений для каждого чата. \\ \hline
 \leftrow 5 & WebSocket-соединение для чата & \leftrow Установка постоянного двунаправленного соединения с сервером через Django Channels; проверка участия пользователя в чате при подключении; поддержка событий \texttt{new\_message} и \texttt{messages\_read}. \\ \hline
 \leftrow 6 & Создание уведомлений на сайте & \leftrow Автоматическое создание записи уведомления в базе данных и немедленная доставка уведомления авторизованному пользователю через WebSocket при наступлении события: новая заявка на бронирование, подтверждение бронирования, отмена бронирования, подтверждение оплаты, новое сообщение, новый отзыв, истечение срока аренды. \\ \hline
 \leftrow 7 & Получение и отметка уведомлений & \leftrow Получение списка уведомлений текущего пользователя; отметка отдельного или всех уведомлений как прочитанных; получение количества непрочитанных уведомлений. \\ \hline
 \leftrow 8 & WebSocket-соединение для уведомлений & \leftrow Установка персонального WebSocket-канала уведомлений для каждого авторизованного пользователя; доставка события \texttt{notification\_created} при появлении нового уведомления. \\ \hline
 \leftrow 9 & Email-уведомления & \leftrow Отправка электронных писем пользователям при ключевых событиях: создание заявки на бронирование, подтверждение, отмена, подтверждение оплаты, новое сообщение в чате, новый отзыв, напоминание о возврате. Отдельная обработка писем с кодом сброса пароля, уведомлений администраторам о заявках на модерацию. \\ \hline
 \leftrow 10 & Напоминания о возврате & \leftrow Управляющая команда Django для выборки подтверждённых бронирований с датой окончания на следующий день и отправки соответствующих email- и сайт-уведомлений арендаторам. \\ \hline
 \leftrow 11 & Генерация договора аренды & \leftrow Автоматическое создание договора аренды при первом обращении к нему по подтверждённому бронированию: генерация уникального номера документа и формирование файла договора в формате PDF с использованием библиотеки ReportLab. \\ \hline
 \leftrow 12 & Содержание PDF-договора & \leftrow Договор содержит: сведения об арендодателе и арендаторе (с учётом типа аккаунта: физлицо, ИП, юрлицо), предмет и срок аренды, арендную плату, права и обязанности сторон, реквизиты, блоки электронной подписи. \\ \hline
 \leftrow 13 & Демонстрационная электронная подпись & \leftrow Подписание договора каждым участником сделки путём ввода имени подписанта и PIN-кода; генерация уникального кода подписи; фиксация даты подписания; отметка договора как полностью подписанного после подписания обеими сторонами; перегенерация PDF с актуальными блоками ЭЦП. \\ \hline
 \leftrow 14 & Справочный раздел FAQ & \leftrow Предоставление структурированного набора часто задаваемых вопросов с пошаговыми инструкциями по основным действиям в системе: создание объявления, заполнение профиля, оформление аренды, оплата, работа с чатом, удаление объявления. \\ \hline
\end{xltabular}

\subsubsection{Требования к API модуля}
Модуль реализован в виде трёх самостоятельных Django-приложений и предоставляет REST API и WebSocket-интерфейсы, используемые клиентской частью приложения.
Модуль предоставляет следующие конечные точки REST API и WebSocket-маршруты.
Подсистема чатов предоставляет следующие REST-endpoint:
\begin{itemize}
    \item \texttt{GET /api/chats/conversations/} — получение списка чатов текущего пользователя с данными о собеседнике, последнем сообщении и счётчике непрочитанных;
    \item \texttt{POST /api/chats/conversations/} — создание нового чата с указанным пользователем и опциональной привязкой к объявлению; при наличии существующего чата возвращается он;
    \item \texttt{GET /api/chats/conversations/\{id\}/} — получение содержимого чата со списком сообщений; при открытии непрочитанные сообщения отмечаются прочитанными;
    \item \texttt{POST /api/chats/conversations/\{id\}/messages/} — отправка нового сообщения (текст и/или изображение) в чат с немедленной рассылкой через WebSocket и созданием уведомления получателю;
    \item \texttt{POST /api/chats/conversations/\{id\}/mark\_read/} — явная отметка сообщений чата как прочитанных;
    \item \texttt{GET /api/chats/support-faq/} — получение структурированного набора вопросов и ответов справочного раздела.
\end{itemize}

Подсистема чатов предоставляет следующий WebSocket-маршрут:
\begin{itemize}
    \item \texttt{ws/chat/\{chat\_id\}/} — персональное подключение к чату; сервер проверяет принадлежность пользователя чату; поддерживаются события \texttt{new\_message} и \texttt{messages\_read}.
\end{itemize}

Подсистема уведомлений предоставляет следующие REST-endpoint:
\begin{itemize}
    \item \texttt{GET /api/notifications/} — получение списка уведомлений текущего пользователя в порядке убывания даты создания;
    \item \texttt{GET /api/notifications/unread\_count/} — получение количества непрочитанных уведомлений;
    \item \texttt{POST /api/notifications/mark\_all\_read/} — отметка всех уведомлений пользователя как прочитанных;
    \item \texttt{POST /api/notifications/\{id\}/mark\_read/} — отметка одного уведомления как прочитанного.
\end{itemize}

Подсистема уведомлений предоставляет следующий WebSocket-маршрут:
\begin{itemize}
    \item \texttt{ws/notifications/} — персональный канал уведомлений; при появлении нового уведомления клиенту доставляется событие \texttt{notification\_created} с полными данными объекта.
\end{itemize}

Подсистема договоров предоставляет следующие REST-endpoint:
\begin{itemize}
    \item \texttt{GET /api/contracts/by-booking/\{booking\_id\}/} — получение или автоматическое создание договора по идентификатору подтверждённого бронирования; ответ содержит текстовый предпросмотр, ссылку на PDF-файл, статус подписания и роль текущего пользователя;
    \item \texttt{POST /api/contracts/\{id\}/sign/} — подписание договора текущим пользователем с передачей имени подписанта и PIN-кода; после подписания обеими сторонами PDF-файл перегенерируется с актуальными блоками ЭЦП.
\end{itemize}

\subsubsection{Требования к моделям данных}

Подсистема чатов использует следующие модели:
\begin{itemize}
    \item \texttt{Chat} — диалог между двумя пользователями с опциональной привязкой к объявлению (поля: участники, объявление, дата создания);
    \item \texttt{Message} — сообщение в чате (поля: чат, отправитель, текст, изображение, дата создания, признак прочтения).
\end{itemize}

Подсистема уведомлений использует следующую модель:
\begin{itemize}
    \item \texttt{Notification} — уведомление пользователя (поля: получатель, тип события, текст сообщения, метаданные в формате JSON для навигации на клиенте, признак прочтения, дата создания). Поддерживаются типы: \texttt{booking\_created}, \texttt{booking\_confirmed}, \texttt{booking\_cancelled}, \texttt{payment\_confirmed}, \texttt{return\_reminder}, \texttt{new\_message}, \texttt{new\_review}, а также события модерации объявлений.
\end{itemize}

Подсистема договоров использует следующую модель:
\begin{itemize}
    \item \texttt{Contract} — договор аренды (поля: бронирование, уникальный номер документа, PDF-файл, данные и даты подписи арендодателя и арендатора, признак полного подписания, дата подписания).
\end{itemize}

\subsubsection{Требования к составу и содержанию PDF-договора}

Формируемый договор аренды в формате PDF содержит следующие разделы:
\begin{itemize}
    \item заголовок с уникальным номером договора и датой его формирования;
    \item преамбулу с описанием сторон с учётом организационно-правовой формы (физическое лицо, индивидуальный предприниматель, юридическое лицо);
    \item раздел «Предмет договора» с наименованием арендуемой вещи и стоимостью аренды;
    \item раздел «Срок действия договора» с датами начала и окончания аренды;
    \item раздел «Права и обязанности сторон»;
    \item раздел «Порядок передачи имущества»;
    \item раздел «Арендная плата и порядок расчётов»;
    \item разделы «Изменение и прекращение договора», «Ответственность сторон», «Разрешение споров», «Прочие условия»;
    \item реквизиты сторон в табличной форме;
    \item блоки демонстрационной электронной подписи арендодателя и арендатора с отображением имени подписанта, даты подписания и кода подписи.
\end{itemize}

\subsection{Моделирование вариантов использования}

Для разрабатываемого модуля была реализована модель, обеспечивающая наглядное представление вариантов использования программы с помощью унифицированного языка визуального моделирования UML. На основании анализа предметной области в программе должны быть реализованы следующие прецеденты:
\begin{enumerate}
    \item Формирование и подписание договора аренды участниками сделки.
    \item Чат между пользователями с возможностью отправить изображение. 
    \item Помощь на часто задаваемые вопросы.
\end{enumerate}

На рисунке \ref{diaprec:image} изображены прецеденты для пользователя.

\begin{figure}[H]
\center{\includegraphics[width=1\linewidth]{diaprec}}
\caption{Диаграмма прецедентов для пользователя}
\label{diaprec:image}
\end{figure}
\subsubsection{Формирование и подписание договора аренды участниками сделки}
Заинтересованные лица и их требования: арендатор и владелец желают сформировать и подписать договор аренды по бронированию.

Предусловие: бронирование подтверждено или завершено.

Постусловие: договор сформирован, после подписи обеих сторон отмечен как подписанный.

Основной успешный сценарий:
\begin{enumerate}
    \item Участник сделки открывает договор по бронированию.
    \item Приложение отправляет запрос на сервер.
    \item Сервер проверяет, что пользователь является участником сделки.
    \item Сервер проверяет статус бронирования.
    \item Сервер создает договор, если он еще не был создан.
    \item Сервер формирует номер договора и файл с данными аренды.
    \item Пользователь вводит имя подписанта и PIN демонстрационной подписи.
    \item Сервер сохраняет подпись участника.
    \item Второй участник выполняет подписание аналогичным образом.
    \item После подписания обеими сторонами сервер отмечает договор как подписанный.
\end{enumerate}

\subsubsection{Чат между пользователями с возможностью отправить изображение}

Заинтересованные лица и их требования: арендатор и владелец желают обсудить детали аренды во встроенном чате.

Предусловие: пользователь авторизован в системе.

Постусловие: сообщение отправлено, сохранено и доступно участникам чата.

Основной успешный сценарий:
\begin{enumerate}
    \item Пользователь открывает карточку вещи или раздел чатов.
    \item Пользователь выбирает собеседника и начинает диалог.
    \item Приложение передает на сервер идентификатор собеседника и при необходимости идентификатор вещи.
    \item Сервер проверяет наличие существующего чата между пользователями.
    \item Если чат отсутствует, сервер создает новый чат.
    \item Пользователь вводит текст сообщения или прикрепляет изображение.
    \item Сервер сохраняет сообщение в базе данных.
    \item Сервер отправляет уведомление получателю.
    \item Приложение отображает новое сообщение в чате.
\end{enumerate}

\subsubsection{Помощь на часто задаваемые вопросы}

Заинтересованные лица и их требования: пользователь желает получить справочную информацию по основным действиям в системе аренды вещей: созданию объявления, заполнению профиля, оформлению аренды, оплате, переписке и просмотру чатов.

Предусловие: пользователь находится на сайте и открыл раздел помощи или окно поддержки.

Постусловие: пользователь получает список часто задаваемых вопросов и ответ с пошаговой инструкцией по выбранной теме.

Основной успешный сценарий:

\begin{enumerate}
    \item Пользователь открывает раздел помощи или окно поддержки на сайте.
    \item Приложение отправляет запрос на сервер для получения списка часто задаваемых вопросов.
    \item Сервер получает подготовленный список FAQ-вопросов из системы.
    \item Сервер передает в приложение список вопросов и связанных с ними ответов.
    \item Приложение отображает пользователю доступные темы помощи.
    \item Пользователь выбирает интересующий вопрос, например о создании объявления, заполнении профиля, бронировании, оплате или чатах.
    \item Приложение отображает ответ по выбранной теме.
    \item Пользователь просматривает пошаговую инструкцию.
    \item При необходимости пользователь переходит в соответствующий раздел сайта и выполняет действие по инструкции.
\end{enumerate}

\subsection{Требования к оформлению документации}

Разработка программной документации и программного изделия должна осуществляться в соответствии с требованиями ГОСТ 19.102-77 и ГОСТ 34.601-90. Единая система программной документации.